% Generated draft: requires mathematical and source review.
\documentclass[11pt]{article}
\usepackage[margin=25mm]{geometry}
\usepackage{fontspec}
\setmainfont{Latin Modern Roman}
\usepackage{amsmath,amssymb,mathtools,unicode-math}
\setmathfont{Latin Modern Math}
\usepackage{enumitem,needspace,xurl,hyperref}
\setlength{\parindent}{0pt}
\setlength{\parskip}{5pt}
\setlength{\emergencystretch}{3em}
\newcommand{\N}{\mathbb{N}}
\newcommand{\Z}{\mathbb{Z}}
\newcommand{\Q}{\mathbb{Q}}
\newcommand{\R}{\mathbb{R}}
\newcommand{\C}{\mathbb{C}}
\newcommand{\F}{\mathbb{F}}
\newcommand{\A}{\mathbb{A}}
\newcommand{\PP}{\mathbb{P}}
\newcommand{\E}{\mathbb{E}}
\newcommand{\eps}{\varepsilon}
\newcommand{\dd}{\,\mathrm{d}}
\begin{document}
\section*{UnsolvedMath ID: 30000497}
\textbf{Source problem code:} catalogue entry
\subsection{Definitions and mathematical statement}
Let $p$ be a prime and let $G$ be a compact $p$-valued group with no elements of order $p$, in the sense intended by the report. The original extraction places the question in the noncommutative case. Write
\[
 \Lambda_G=\Z_p[[G]]
    =\varprojlim_{U}\Z_p[G/U],
\]
where $U$ ranges over open normal subgroups; this is the usual integral Iwasawa-algebra convention. Let $M$ be a finitely generated torsion $\Lambda_G$-module, as specified in the restored statement.

The questions are:
\begin{enumerate}
\item Can the reflexive right ideals $L_i$ in the report's structure theorem be replaced, possibly after changing the choices allowed by that theorem, by principal right ideals? A right ideal is principal if $L=a\Lambda_G$ for some $a\in\Lambda_G$. Reflexivity means that the canonical evaluation map to the module's double dual is an isomorphism, with the appropriate right/left module duals over $\Lambda_G$. The requested replacement must preserve the structure-theorem description of $M$; it is not the assertion that every reflexive right ideal in the ring is principal.
\item How should a characteristic element be defined in general for the intended torsion modules over this noncommutative Iwasawa algebra?
\end{enumerate}

The structure theorem itself is not supplied: the role of the ideals, the side of the module $M$, and the permitted equivalence or modification of representatives are missing. The original wording incorrectly or differently refers to a theorem ``for pseudonull modules,'' whereas the supplied verification expressly restores a finitely generated torsion-module hypothesis. This discrepancy is retained rather than treating torsion and pseudonull as synonymous. Torsion requires the report's rank or localization convention; no unprovided Ore-localization hypotheses are invented.

For the characteristic-element question the target group, existence properties, ambiguity up to units or another equivalence, and functorial or exact-sequence requirements are not specified. The dated assessment reports localized $K_1$ constructions for broad classes, but does not state a localization or boundary map covering every module here. Those constructions do not by themselves answer the stronger principal-ideal question or define an ordinary ring element in complete generality.
\subsection{Short English statement}
The first question asks whether the ideals appearing in a torsion-module structure theorem can be chosen principal while preserving that description. The second asks for a general noncommutative characteristic element. The exact theorem, allowed replacements, and characteristic-element target are missing.
\subsection{Sources}
\begin{itemize}
\item Local UnsolvedMath registry, permanent ID 30000497, code catalogue entry, ``Principal Ideals and Characteristic Elements for Iwasawa Modules.'' Proposed year 2006; proposer unspecified.
\item \emph{Pro-$p$ Extensions of Global Fields and pro-$p$ Groups} (2007), Oberwolfach Reports, pp.\ 1463--1536, DOI 10.4171/owr/2006/25; \url{https://doi.org/10.4171/owr/2006/25}. The displayed report year 2007 differs from the DOI and proposed year 2006; escaped dollar signs in the title are normalized as mathematical $p$.
\item Supplied ``Venjakob characteristic elements'': \url{https://doi.org/10.1515/crll.2005.2005.583.193}.
\item Supplied ``recent noncommutative characteristic-element work'': \url{https://arxiv.org/abs/2412.13812}.
\item Statement status corrected\_verified; the supplied verification says the official report was checked directly and torsion-module, reflexive-ideal, and characteristic-element clauses restored.
\item Literature triage checked 2026-08-21: partially solved. It reports localized $K_1$ characteristic elements in broad classes, while leaving the principal-reflexive-ideal formulation and completely general construction unresolved. Exact hypotheses and current status require independent verification. Created 2026-08-25; updated 2026-09-16.
\item Registry curation by Ulam AI / UnsolvedMath Contributors, CC BY 4.0; underlying report and literature provenance remain as supplied. Unreviewed statement-only draft.
\end{itemize}
\end{document}
